\documentclass[11pt,letterpaper]{article}

\usepackage[T1]{fontenc}
\usepackage{amsmath,amssymb}
\usepackage{tikz}
\usepackage{pgfplots}
\pgfplotsset{compat=1.18}
\usepackage{algorithm}
\usepackage{algpseudocode}
\usepackage{booktabs}
\usepackage{enumitem}
\usepackage[hidelinks]{hyperref}

\newcommand{\N}{\mathbb{N}}
\newcommand{\Z}{\mathbb{Z}}
\newcommand{\Q}{\mathbb{Q}}
\newcommand{\R}{\mathbb{R}}
\newcommand{\C}{\mathbb{C}}
\newcommand{\Map}[3]{#1 \colon #2 \to #3}
\newcommand{\Mod}[1]{#1\text{-}\mathrm{Mod}}
\newcommand{\GroupHom}{\mathrm{Hom}_{\mathbf{Grp}}}
\newcommand{\RingHom}{\mathrm{Hom}_{\mathbf{Ring}}}

\newtheorem{theorem}{Theorem}
\newtheorem{lemma}[theorem]{Lemma}
\newtheorem{corollary}[theorem]{Corollary}
\newtheorem{definition}[theorem]{Definition}
\newtheorem{remark}[theorem]{Remark}
\newenvironment{proof}{\noindent\textit{Proof.}}{\hfill\(\square\)}

\title{Immediate return at exponent zero}
\author{}
\date{\today}

\begin{document}
\maketitle

\begin{abstract}
At time zero the affine state of an arbitrary integer seed is the seed itself.
The requirement that the signed divisor divide no earlier affine state is a
bounded quantifier over the empty initial segment of the natural numbers, hence
is true for every predicate. A first hit at time zero is therefore exactly
divisibility of the seed. On that event the raw step is division by the signed
divisor, absolute values multiply back to the absolute value of the seed, and
the core comparison holds with slope \(1\) and constant \(0\). The argument uses
\(|d|\ge 2\) only to obtain the strict inequality \(1<|d|\). No positive
exponent is treated.
\end{abstract}

\section{Definitions}

Throughout, a parameter record \(p\) consists of integers \(q\), \(c\), and
\(d\) with \(|d|\ge 2\). In particular \(d\neq 0\). The raw step is
\[
T(n)=\begin{cases}
n/d, & d\mid n,\\
qn+c, & d\nmid n.
\end{cases}
\]
The offset sum is defined by \(U_0(q)=0\) and
\(U_{j+1}(q)=U_j(q)+q^j\). Thus \(U_0(q)\) is the empty sum.

\begin{definition}
\label{def:affine-state}
For \(q,c,n\in\Z\) and \(j\in\N\), the affine state is
\[
A_j(q,c,n)=q^j n+c\,U_j(q).
\]
\end{definition}

\begin{definition}
\label{def:first-hit}
The time \(j\in\N\) is a first hit of the seed \(n\in\Z\), written
\(\operatorname{Hit}(p,n,j)\), when
\[
d\mid A_j(q,c,n)
\]
and, for every natural number \(i<j\),
\[
d\nmid A_i(q,c,n).
\]
The second clause is the bounded quantifier \(\forall i\,(i<j\rightarrow
d\nmid A_i(q,c,n))\).
\end{definition}

\section{The empty initial segment}

\begin{lemma}
\label{lem:affine-zero}
For all \(q,c,n\in\Z\),
\[
A_0(q,c,n)=n.
\]
\end{lemma}

\begin{proof}
By Definition~\ref{def:affine-state} and the clauses \(q^0=1\) and \(U_0(q)=0\),
\[
A_0(q,c,n)=q^0 n+c\,U_0(q)=1\cdot n+c\cdot 0.
\]
The integer identities \(1\cdot n=n\), \(c\cdot 0=0\), and \(n+0=n\) give
\(A_0(q,c,n)=n\).
\end{proof}

\begin{lemma}
\label{lem:no-earlier}
Let \(P\) be any predicate on \(\N\). Then \(\forall i\in\N,\ i<0\rightarrow P(i)\).
\end{lemma}

\begin{proof}
By definition, the bounded quantifier over the initial segment of natural
numbers strictly less than \(0\) is the ordinary quantifier
\[
\forall i\in\N,\quad (i<0\rightarrow P(i)).
\]
Fix \(i\in\N\). Every natural number is at least \(0\), so the statement
\(i<0\) is false. A false hypothesis implies \(P(i)\), whatever the predicate
\(P\) is. The range is empty: the implication is not obtained by restricting
\(P\), nor by ignoring the quantifier.
\end{proof}

\begin{theorem}
\label{thm:first-hit-zero}
For every parameter record \(p\) with divisor \(d\), and every \(n\in\Z\),
\[
\operatorname{Hit}(p,n,0)\quad\text{if and only if}\quad d\mid n.
\]
\end{theorem}

\begin{proof}
Unfold Definition~\ref{def:first-hit} at time \(0\). The statement
\(\operatorname{Hit}(p,n,0)\) is the conjunction of \(d\mid A_0(q,c,n)\) and
\[
\forall i\in\N,\quad i<0\rightarrow d\nmid A_i(q,c,n).
\]
Lemma~\ref{lem:affine-zero} gives \(A_0(q,c,n)=n\), so the first conjunct is
equivalent to \(d\mid n\).

The second conjunct is Lemma~\ref{lem:no-earlier} applied to the predicate
\(i\mapsto d\nmid A_i(q,c,n)\). It is therefore true, and it does not depend
on \(n\) or on \(d\).

If \(\operatorname{Hit}(p,n,0)\) holds, the first conjunct and
Lemma~\ref{lem:affine-zero} yield \(d\mid n\). Conversely, if \(d\mid n\), then
\(d\mid A_0(q,c,n)\) by the same lemma, and the second conjunct holds by
Lemma~\ref{lem:no-earlier}. Both conjuncts are proved; the earlier-time
conjunct is not omitted in the converse direction.
\end{proof}

\section{The immediate quotient}

\begin{theorem}
\label{thm:immediate-quotient}
Let \(p\) be a parameter record with divisor \(d\), and let \(n\in\Z\) satisfy
\(d\mid n\). Then the raw step equals the quotient,
\[
T(n)=n/d,
\]
and the absolute values cancel,
\[
|d|\cdot |n/d|=|n|.
\]
\end{theorem}

\begin{proof}
The hypothesis \(d\mid n\) selects the division branch of the raw step, so
\(T(n)=n/d\) by the definition of \(T\). The same hypothesis is the integer
cancellation law
\[
d\cdot(n/d)=n.
\]
Absolute value is multiplicative, hence
\[
|d\cdot(n/d)|=|d|\cdot|n/d|.
\]
The left-hand side equals \(|n|\) because \(d\cdot(n/d)=n\). Therefore
\(|d|\cdot|n/d|=|n|\).
\end{proof}

\begin{theorem}
\label{thm:immediate-core}
Let \(p\) be a parameter record with divisor \(d\), and let \(n\in\Z\) satisfy
\(d\mid n\). Let \(A,B\in\N\) satisfy \(A=1\) and \(B=0\). Then
\[
A<|d|
\]
and
\[
|d|\cdot|n/d|\le A\,|n|+B.
\]
\end{theorem}

\begin{proof}
The record supplies \(|d|\ge 2\). Since \(1<2\) and \(2\le|d|\), transitivity
of the strict and non-strict orders on \(\N\) yields \(1<|d|\). The hypothesis
\(A=1\) rewrites this as \(A<|d|\).

Theorem~\ref{thm:immediate-quotient} gives \(|d|\cdot|n/d|=|n|\). Equality
implies the non-strict comparison \(|d|\cdot|n/d|\le|n|\). The remaining
hypotheses are \(A=1\) and \(B=0\), so
\[
A\,|n|+B=1\cdot|n|+0=|n|,
\]
where \(1\cdot|n|=|n|\) and \(|n|+0=|n|\). Substituting this right-hand side
yields
\[
|d|\cdot|n/d|\le A\,|n|+B,
\]
as claimed. The inequality \(|d|\ge 2\) is used for \(A<|d|\) and is not
required for the absolute-value cancellation.
\end{proof}

\begin{remark}
The two conclusions are the core comparison for the single map \(n\mapsto n/d\)
on the set of integers divisible by \(d\), with slope \(1\) and constant \(0\).
They do not assert the same bound for a return taken after one or more affine
steps, nor do they assert that every seed is eventually divisible by \(d\).
\end{remark}

\end{document}
